\begin{table}[htbp]\centering
\begin{threeparttable}
\caption{Cultures, named species and species complexes that form a single clade (are monophyletic) on trees of
the culture copies.}
\label{tab:monophyly}
\footnotesize\setlength{\tabcolsep}{4pt}
\begin{tabular}{lrrrrrr}
\toprule
 & \multicolumn{2}{c}{Cultures} & \multicolumn{2}{c}{Named species} & \multicolumn{2}{c}{Complexes}\\
\cmidrule(lr){2-3}\cmidrule(lr){4-5}\cmidrule(lr){6-7}
Tree & all & main & all & main & all & main\\
\midrule
SSU flank & 17/139 & 19/138 & 28/37 & 28/37 & 3/4 & 4/4\\
Full ITS & 18/139 & 24/138 & 24/37 & 29/37 & 2/4 & 3/4\\
LSU flank & 19/139 & 22/138 & 28/37 & 30/37 & 3/4 & 4/4\\
Linked copy & 23/139 & 28/138 & 31/37 & 35/37 & 3/4 & 4/4\\
\bottomrule
\end{tabular}
\begin{tablenotes}\footnotesize
\item Monophyletic groups / groups tested; a group is tested when it has two or more copies. \emph{all}: all copies;
\emph{main}: the same trees with the divergent copies removed. One tree per order and region (maximum likelihood,
IQ-TREE); the linked copy joins the three regions. Alignments and trees are in Supplementary Data S1.
\end{tablenotes}
\end{threeparttable}
\end{table}
